\documentclass[12pt,a4paper]{article}

% encoding / language
\usepackage[utf8]{inputenc}
\usepackage[T1]{fontenc}
\usepackage[brazilian]{babel}

% misc
\usepackage{mathtools}
\usepackage{enumitem}
\usepackage{fancyvrb}
\usepackage[dvipsnames,table]{xcolor}
\usepackage{float}

% citations
\usepackage{csquotes}
\usepackage[style=abnt,maxbibnames=5]{biblatex}
\addbibresource{ref.bib}

% spacing
\usepackage{indentfirst}
\frenchspacing
\setlength{\parindent}{1.25cm}
\usepackage{setspace}
\onehalfspacing
\usepackage[top=3cm,bottom=2cm,left=3cm,right=2cm]{geometry}

% font
\usepackage[scaled]{helvet}
\renewcommand{\familydefault}{\sfdefault}
\usepackage{sfmath}

% hypertext & pdf
\usepackage{hyperref}
\hypersetup{
    colorlinks=true,
    urlcolor=[RGB]{0, 0, 130},
    linkcolor=.,
    citecolor=.,
    pdftitle={Enade SO}
}
\urlstyle{same}

% graphics
\usepackage{graphicx}
\graphicspath{{graphics}}

\newenvironment{Question}[1]{%
    \pdfbookmark[3]{Questão #1}{q:#1}
    \label{q:#1}
    \par\bigskip\noindent
    \textbf{Questão #1} \dotfill
    \par\noindent
    \nopagebreak
    \ignorespaces
}{%
    \noindent\dotfill
}

\newenvironment{QuestionOptions}{%
    \begin{spacing}{0.9}%
    \small
    \begin{enumerate}[label=\Alph*)]
}{%
    \end{enumerate}
    \end{spacing}
}

\newcommand{\QuestionFigure}[2]{
    \begin{figure}[H]
    \centering
    \includegraphics[alt={{#2}},width=\linewidth]{#1}
    \end{figure}
}

\newenvironment{QuestionPage}[1][\linewidth]{%
    \medskip\noindent
    \begin{minipage}[t]{#1}
}{%
    \end{minipage}
    \medskip
}

\newcommand{\unit}[1]{\,\mathrm{#1}}

\renewcommand{\theFancyVerbLine}{\ttfamily\color{gray}\arabic{FancyVerbLine}}

\title{\color{red}Rascunho}
\author{
    \small\href{https://github.com/eemidioo/enade-so}{Código-fonte GitHub}%
    \thanks{\url{https://github.com/eemidioo/enade-so}}
}
\date{\today}

\begin{document}
\maketitle
2ª atividade da disciplina optativa ENADE-2026.2, do Professor Ricardo Azevedo Moreira da Silva.
Este documento é destinado a estudantes de Ciência da Computação
e contém um resumo de conceitos fundamentais de Sistemas Operacionais,
baseado nas questões objetivas selecionadas das três provas de 2014, 2017 e 2021 do \citefield{enade}{title}.
Ressalta-se que a reprodução das questões nem sempre é \textit{ipsis litteris}, ou seja,
algumas adaptações são feitas para deixar o texto mais conciso ou consistente com o conteúdo apresentado.

\section{Introdução}
O Sistema Operacional (SO) é o software fundamental que gerencia o hardware do computador e fornece uma base para a execução de programas.
Ele atua como um gerenciador de recursos, distribuindo tempo de processador, memória e dispositivos de entrada e saída
de forma segura e eficiente \parencite[cap.~1]{tanenbaum2016}.

Quando um computador executa múltiplas tarefas de forma intercalada,
dizemos que há multiprogramação e que estamos lidando com \emph{sistemas concorrentes}.
Nesse cenário, o SO gerencia a ilusão de que vários processos estão rodando ao mesmo tempo alternando rapidamente a CPU entre eles.
Com o advento dos processadores de múltiplos núcleos, alcançamos também o \emph{processamento paralelo},
onde instruções de diferentes processos ou threads são executadas literalmente no mesmo instante de tempo físico,
em núcleos físicos distintos.
O gerenciamento adequado da concorrência e do paralelismo é um dos pilares modernos
do design de sistemas operacionais \parencite[cap.~2]{tanenbaum2016}.

\section{Processos}
Um processo é, em essência, um programa em execução.
Como os recursos do computador são limitados e múltiplos processos competem pela CPU,
o sistema operacional utiliza um \emph{escalonador} para decidir qual processo deve ser executado a seguir
e por quanto tempo.
Diferentes ambientes exigem diferentes estratégias de escalonamento \parencite[cap.~2]{tanenbaum2016}.

Para que esses processos inicialmente isolados possam cooperar,
utiliza-se a \emph{comunicação entre processos} (\textit{Interprocess Communication} - IPC).
O IPC fornece recursos como memória compartilhada, \textit{pipes} e troca de mensagens
para viabilizar a transferência de dados e a sincronização entre eles sem violar a proteção do sistema \parencite[seç.~2.3]{tanenbaum2016}.

\subsection{Exemplos}
Seguem agora questões resolvidas do Enade que envolvem os processos em um sistema operacional.

\begin{Question}{2021-09}
Quando um computador é multiprogramado, ele geralmente tem múltiplos processos ou \textit{threads} que competem pela CPU ao mesmo tempo.
Essa situação ocorre sempre que dois ou mais processos estão simultaneamente no estado pronto.
Se somente uma CPU se encontrar disponível, deverá ser feita uma escolha de qual processo executar em seguida.
A parte do sistema operacional que faz a escolha é chamada de \emph{escalonador},
e o algoritmo que ele usa é o \emph{algoritmo de escalonamento}.
\vspace{-0.5em}
\begin{flushright}
    \scriptsize
    TANENBAUM, A. S. \textbf{Sistemas Operacionais Modernos}.\\
    3. ed., São Paulo: Pearson, 2010 (adaptado).
\end{flushright}
Considerando que em ambientes diferentes são necessários algoritmos diferentes de escalonamento,
garantindo assim que seja maximizado o uso de seus recursos,
assinale a opção que apresenta um algoritmo de escalonamento seguido do tipo de ambiente no qual deva ser implementado.
\begin{QuestionOptions}
    \item Primeiro a chegar, último a sair (\textit{first in, last out} - FILO); propício para sistemas de tempo real.
    \item Escalonamento por taxas monotônicas (\textit{rate monotonic scheduling} - RMS); propício para sistemas em lote.
    \item Tarefa mais curta primeiro; propício para sistemas interativos.
    \item Escalonamento por chave circular (\textit{round-robin}); propício para sistemas de tempo real.
    \item Escalonamento por prioridades; propício para sistemas interativos.
\end{QuestionOptions}
\end{Question}

Para resolver essa questão, devemos relembrar as categorias de escalonamento citadas por \textcite[seç.~2.4]{tanenbaum2016}.
Em sistemas em lote (\textit{batch}), onde não há usuários esperando resposta imediata,
algoritmos como ``Primeiro a chegar, primeiro a ser servido (FIFO)'' e ``Tarefa mais curta primeiro (SJF)'' são adequados.
Em sistemas interativos, onde o tempo de resposta é crucial, utilizam-se o \textit{Round-Robin} (chave circular)
e o Escalonamento por Prioridades.
Em sistemas de tempo real, o foco principal é o cumprimento rigoroso de prazos (\textit{deadlines}),
utilizando-se algoritmos estáticos (tomam decisões antes do sistema ser executado)
ou dinâmicos (tomam decisões em tempo de execução) específicos para esse fim.

Analisando as opções:
a opção A cita FILO, o que seria uma pilha, inadequado para escalonamento justo;
a opção B inverte RMS (tempo real) para sistemas em lote;
a opção C sugere tarefa mais curta para sistemas interativos (na verdade, é para lote);
a opção D sugere \textit{Round-Robin} para tempo real (o que não garante cumprimento de prazos rígidos).
A opção correta é a E, pois o escalonamento por prioridades permite que tarefas interativas importantes tenham a resposta rápida necessária.

\section{Gerenciamento de memória}
O gerenciamento de memória é responsável por alocar o espaço necessário para os processos
e proteger a memória de um processo contra interferências de outros.
A técnica moderna mais utilizada é a \emph{paginação}, que divide a memória lógica em blocos de tamanho fixo (páginas)
e a memória física em blocos equivalentes (quadros ou \textit{frames}),
resolvendo problemas graves como a fragmentação externa \parencite[cap.~3]{tanenbaum2016}.

A fragmentação interna ocorre quando a memória é alocada em partições ou páginas de tamanho fixo $S_{\text{bloco}}$,
de modo que a requisição de um processo $S_{\text{processo}} < S_{\text{bloco}}$
resulta em um desperdício interno $D = S_{\text{bloco}} - S_{\text{processo}}$ dentro do próprio bloco atribuído,
o qual não pode ser aproveitado por outros programas;
já a fragmentação externa surge na alocação dinâmica de tamanho variável
quando o total de memória livre disponível é $M \ge S_{\text{processo}}$,
contudo esse espaço encontra-se disperso em pequenos segmentos não contíguos entre os blocos ocupados,
impedindo o carregamento contínuo do novo processo \parencite[cap.~3]{tanenbaum2016}.

\subsection{Exemplos}
Seguem agora questões resolvidas do Enade que envolvem o gerenciamento de memória.

\begin{Question}{2014-11}
Considere um esquema de gerência de memória por paginação simples, onde a memória física é dividida
em quadros (\textit{frames}) de $1 \unit{Kbyte}$ e endereçada por byte.
Os espaços de endereçamento dos processos são múltiplos de $1 \unit{Kbyte}$.
A tabela de páginas para um determinado processo P é apresentada a seguir,
em que o primeiro bit (BV) mostra se a página é válida (1) ou inválida (0).
\begin{center}
\begin{tabular}{|c|c|c|}
\hline
\textbf{Página} & \textbf{BV} & \textbf{Quadro (\textit{frame})} \\
\hline
0 & 1 & 0010 \\
1 & 1 & 0100 \\
2 & 1 & 0001 \\
3 & 1 & 0111 \\
4 & 1 & 0000 \\
5 & 1 & 1101 \\
6 & 0 & 1111 \\
7 & 0 & 0110 \\
\hline
\end{tabular}
\end{center}
Com base na tabela apresentada, avalie as afirmações a seguir.
\begin{enumerate}[label=\Roman*.]
    \item O endereço físico é composto por 13 bits.
    \item O esquema de gerência de memória apresentado reduz a fragmentação externa.
    \item A tradução do endereço lógico \texttt{01100000000110} para endereço físico causa exceção.
\end{enumerate}
É correto o que se afirma em
\begin{QuestionOptions}
    \item I, apenas.
    \item II, apenas.
    \item I e III, apenas.
    \item II e III, apenas.
    \item I, II e III.
\end{QuestionOptions}
\end{Question}

A afirmação I é falsa.
O tamanho de cada quadro é de $1 \unit{Kbyte}$ (considerando $\unit{K} = 1024$),
o que requer 10 bits de deslocamento ($2^{10} = 1024 \unit{bytes}$).
Observando a tabela, o número do quadro utiliza 4 bits (de \texttt{0000} a \texttt{1111}).
Portanto, o endereço físico completo é a junção dos 4 bits do quadro com os 10 bits de deslocamento,
totalizando 14 bits, e não 13.

A afirmação II é verdadeira.
A paginação divide a memória em blocos de tamanho fixo,
permitindo que qualquer quadro livre receba uma página de qualquer processo.
Isso elimina a fragmentação externa,
restando apenas a possibilidade de fragmentação interna na última página do processo \parencite[seç.~3.3, seç.~3.5.3, p.~196]{tanenbaum2016}.

A afirmação III é falsa.
Como a tabela possui 8 páginas (índices de 0 a 7),
são necessários 3 bits para identificar a página virtual ($2^3 = 8$).
Dado que os 10 bits menos significativos do endereço lógico correspondem ao deslocamento dentro da página,
os 3 bits mais significativos indicam o número da página.
Nos bits mais significativos de \texttt{01100000000110}, temos \texttt{011},
o que equivale à página 3 em decimal.
Consultando a tabela para a página 3,
constata-se que o bit de validade (BV) é 1 e o quadro correspondente é \texttt{0111}.
Como o bit de validade é igual a 1,
a página está presente na memória física (RAM) e o endereço é traduzido normalmente para o quadro \texttt{0111},
sem gerar exceção ou falta de página (\textit{page fault}) \parencite[seç.~3.3]{tanenbaum2016}.
A opção correta é a B.

\begin{Question}{2017-29}
O projetista do gerenciador de memória de um novo sistema operacional precisa escolher
entre os algoritmos de substituição de páginas FIFO (\textit{First In First Out}) e LRU (\textit{Least Recently Used}).
Para isso, avaliou o número de faltas de página obtidas em ambos os algoritmos para o tamanho de memória de 4 páginas,
utilizando a sequência de acessos às páginas \mbox{1-2-3-4-1-2-5-1-2-3-4-5} de um processo e memória inicialmente vazia.\\\\
Com base nessa simulação, avalie as asserções a seguir e a relação proposta entre elas.
\begin{enumerate}[label=\Roman*.]
    \item Na simulação proposta, é possível observar que os algoritmos FIFO e LRU apresentam o mesmo desempenho.
\end{enumerate}
\begin{center}\textbf{PORQUE}\end{center}
\begin{enumerate}[label=\Roman*.,start=2]
    \item Os parâmetros utilizados na simulação são insuficientes para determinar a diferença de funcionamento entre os algoritmos.
\end{enumerate}
A respeito dessas asserções, assinale a opção correta.
\begin{QuestionOptions}
    \item As asserções I e II são proposições verdadeiras, e a II é uma justificativa correta da I.
    \item As asserções I e II são proposições verdadeiras, mas a II não é uma justificativa correta da I.
    \item A asserção I é uma proposição verdadeira, e a II é uma proposição falsa.
    \item A asserção I é uma proposição falsa, e a II é uma proposição verdadeira.
    \item As asserções I e II são proposições falsas.
\end{QuestionOptions}
\end{Question}

Para avaliar as asserções, precisamos simular ambos os algoritmos com 4 quadros e a sequência 1, 2, 3, 4, 1, 2, 5, 1, 2, 3, 4, 5.
No FIFO, a página que está na memória há mais tempo é a escolhida para ser removida:
\begin{footnotesize}
\begin{itemize}
    \item 1, 2, 3, 4: causam 4 faltas (quadros ficam cheios: [1, 2, 3, 4]).
    \item 1, 2: acertos (já estão na memória).
    \item 5: falta, substitui o 1 (quadros: [5, 2, 3, 4]).
    \item 1: falta, substitui o 2 (quadros: [5, 1, 3, 4]).
    \item 2: falta, substitui o 3 (quadros: [5, 1, 2, 4]).
    \item 3: falta, substitui o 4 (quadros: [5, 1, 2, 3]).
    \item 4: falta, substitui o 5 (quadros: [4, 1, 2, 3]).
    \item 5: falta, substitui o 1 (quadros: [4, 5, 2, 3]).
\end{itemize}
\end{footnotesize}
Total no FIFO: 10 faltas de página.
No LRU, a página que não é utilizada há mais tempo é removida:
\begin{footnotesize}
\begin{itemize}
    \item 1, 2, 3, 4: causam 4 faltas (quadros: [1, 2, 3, 4]).
    \item 1, 2: acertos (1 e 2 passam a ser as mais recentemente usadas).
    \item 5: falta, substitui o 3 (quadros: [1, 2, 5, 4]).
    \item 1, 2: acertos (novamente, são as mais recentes).
    \item 3: falta, substitui o 4 (quadros: [1, 2, 5, 3]).
    \item 4: falta, substitui o 5 (quadros: [1, 2, 4, 3]).
    \item 5: falta, substitui o 1 (quadros: [5, 2, 4, 3]).
\end{itemize}
\end{footnotesize}
Total no LRU: 8 faltas de página \parencite[seç.~3.4]{tanenbaum2016}.

A asserção I é falsa, pois os desempenhos foram diferentes (10 faltas no FIFO contra 8 no LRU).
A asserção II também é falsa, uma vez que a sequência de acessos e a quantidade de quadros fornecidos
no enunciado são parâmetros suficientes para realizar as simulações e constatar a diferença.
Portanto, a opção E é a correta.

\section{Concorrência e Threads}
A \emph{concorrência} refere-se à capacidade de gerenciar a execução intercalada de múltiplos fluxos de controle no tempo,
enquanto o \emph{paralelismo} exige o uso de múltiplos núcleos físicos para executar
instruções simultaneamente no mesmo instante \parencite[cap.~2]{tanenbaum2016}.

As \textit{threads} (ou linhas de execução) representam uma das principais formas de materializar a concorrência,
permitindo que múltiplos fluxos de controle compartilhem o mesmo espaço de memória e sejam intercalados ao longo do tempo.
Isso facilita o compartilhamento de dados, mas traz desafios profundos;
pois quando múltiplos processos ou \textit{threads} acessam recursos
ou dados compartilhados de forma concorrente sem os devidos mecanismos necessários,
ocorrem as \emph{condições de corrida} (\textit{race conditions}),
nas quais o resultado final da execução torna-se imprevisível
e dependente da ordem e do momento exato em que o escalonador alterna as instruções \parencite[cap.~2]{tanenbaum2016}.

Para evitar as condições de corrida,
a parte do programa que acessa a memória ou recurso compartilhado é denominada \emph{região crítica} (\textit{critical section}).
A solução para esse problema exige garantir a \emph{exclusão mútua}, isto é, assegurar que,
quando uma \textit{thread} estiver executando instruções dentro de sua região crítica,
nenhuma outra \textit{thread} possa entrar na sua própria região crítica
para acessar os mesmos dados compartilhados simultaneamente \parencite[seç.~2.3]{tanenbaum2016}.

Em um sistema computacional podemos ter vários recursos que podem ser usados somente por uma linha de execução de cada vez.
Dessa forma, um \emph{impasse} (\textit{deadlock}) ocorre quando um conjunto de processos fica permanentemente bloqueado
porque cada um detém um recurso e aguarda a liberação de outro mantido por um processo do mesmo grupo.
A ocorrência de um impasse exige a satisfação simultânea de quatro condições necessárias:
exclusão mútua (um recurso está associado a apenas um processo ou está disponível),
posse e espera (processos que já possuem recursos, mas esperam por mais recursos),
não preempção (recursos que precisam ser explicitamente liberados)
e espera circular (cada processo esperando um recurso na posse de outro processo seguinte da cadeia).
As abordagens para tratar o problema envolvem a ignorância do evento
(quando a ocorrência não é tão significativa;
por exemplo, se acontece uma vez a cada cinco anos, sem grandes prejuízos e com uma remediação suficiente),
a detecção seguida de recuperação, a evitação pela alocação cuidadosa de recursos
e a prevenção pela negação de uma das quatro condições.

Uma trava (\textit{lock}) garante a exclusão mútua ao permitir que apenas uma \textit{thread} por vez acesse um recurso compartilhado.
Já os \emph{semáforos} utilizam um contador manipulado por operações atômicas
(\texttt{down} e \texttt{up}~--- isto é, executadas como uma unidade indivisível de instrução,
sem possibilidade de interrupção ou interferência de outras \textit{threads});
ao contrário dos \textit{locks},
eles podem atuar tanto no modo binário (funcionando como uma trava) quanto de contagem,
permitindo gerenciar o acesso simultâneo a um número limitado de instâncias de um recurso \parencite[cap.~2,6]{tanenbaum2016}.

\subsection{Exemplos}
Seguem agora questões resolvidas do Enade que envolvem \textit{threads}.

\begin{Question}{2014-20}
Um processo tem um ou mais fluxos de execução, normalmente denominados apenas por \textit{threads}.
\QuestionFigure{2014-20}{TODO}
\vspace{-0.5em}
\begin{flushright}
    \scriptsize
    TANENBAUM, A. S. \textbf{Sistemas Operacionais Modernos}.\\
    3. ed., São Paulo: Pearson, 2010 (adaptado).
\end{flushright}
A partir das figuras 1 e 2 apresentadas, avalie as afirmações a seguir.
\begin{enumerate}[label=\Roman*.]
    \item Tanto na figura 1 quanto na figura 2, existem três \textit{threads} que utilizam o mesmo espaço de endereçamento.
    \item Tanto na figura 1 quanto na figura 2, existem três \textit{threads} que utilizam três espaços de endereçamento distintos.
    \item Na figura 2, existe um processo com um único espaço de endereçamento e três \textit{threads} de controle.
    \item Na figura 1, existem três processos tradicionais, cada qual tem seu espaço de endereçamento e uma única \textit{thread} de controle.
    \item As \textit{threads} permitem que várias execuções ocorram no mesmo ambiente de processo de forma independente uma das outras.
\end{enumerate}
É correto apenas o que se afirma em
\begin{QuestionOptions}
    \item I, II e III.
    \item I, II e IV.
    \item I, III e V.
    \item II, IV e V.
    \item III, IV e V.
\end{QuestionOptions}
\end{Question}

A afirmação I é falsa, pois na figura 1 os processos são isolados, não compartilhando a mesma memória.
A afirmação II é falsa, pois na figura 2 as três \textit{threads} dividem o mesmo espaço de endereçamento.
A afirmação III é verdadeira, descrevendo precisamente o modelo de concorrência por \textit{threads} em um processo demonstrado na figura 2.
A afirmação IV é verdadeira, descrevendo o modelo tradicional multiprogramado representado na figura 1.
A afirmação V também é verdadeira e define a premissa conceitual fundamental das \textit{threads} \parencite[seç.~2.2]{tanenbaum2016}.
A opção correta é a E.

\begin{Question}{2014-21}
O fragmento de código a seguir, escrito em Java, descreve duas implementações diferentes para um \textit{lock}.
Ambas possuem um método denominado \texttt{acquire} e um método denominado \texttt{release}.

\begin{QuestionPage}[0.4\linewidth]
\begin{Verbatim}
class LockA {
    private int turn = 0;
    public void acquire(int tid) {
        while (turn == (1 - tid));
    }
    public void release(int tid) {
        turn = (1 - tid);
    }
}
\end{Verbatim}
\end{QuestionPage}
\hfill
\begin{QuestionPage}[0.4\linewidth]
\begin{Verbatim}
class LockB {
    public void acquire() {
        disableInterrupts();
    }
    public void release() {
        enableInterrupts();
    }
}
\end{Verbatim}
\end{QuestionPage}

\noindent
[Adaptado:] Considera-se que as duas implementações são utilizadas por no máximo duas \textit{threads} e o código foi desenvolvido
para ser utilizado em uma máquina com um ou dois processadores.
A partir das informações apresentadas, avalie as afirmações a seguir.
\begin{enumerate}[label=\Roman*.]
    \item A implementação de LockA garante progresso.
    \item A implementação de LockB garante progresso.
    \item A implementação de LockA garante exclusão mútua.
    \item A implementação de LockB garante exclusão mútua.
\end{enumerate}
É correto apenas o que se afirma em
\begin{QuestionOptions}
    \item I e II.
    \item II e III.
    \item III e IV.
    \item I, II e IV.
    \item I, III e IV.
\end{QuestionOptions}
\end{Question}

O LockA implementa o método de alternância explícita.
\textcite[seç.~2.3.2--3]{tanenbaum2016} demonstram que esse método garante a exclusão mútua (Afirmação III verdadeira),
pois o valor da variável \texttt{turn} dita explicitamente quem pode entrar.
No entanto, ela viola o princípio do progresso (Afirmação I falsa).
Se a thread 0 sair da região crítica e atualizar \texttt{turn} para 1,
mas a thread 1 não tiver interesse em entrar na região crítica, a thread 0 ficará bloqueada para sempre na próxima tentativa,
mesmo a região estando livre; causando um impasse (\textit{deadlock}).

O LockB tenta garantir a exclusão mútua desabilitando as interrupções de hardware.
Essa técnica não impede que a thread tome a posse da trava (Afirmação II verdadeira, há progresso),
mas falha na garantia de exclusão mútua em sistemas com mais de um processador (Afirmação IV falsa).
O enunciado ressalta que o sistema pode ter dois processadores.
Desabilitar as interrupções afeta apenas o processador que executou a instrução;
a thread rodando no outro processador pode invadir a memória compartilhada \parencite[seç.~2.3.1]{tanenbaum2016}.
A opção correta é a B.

\begin{Question}{2017-31}
Considere o programa a seguir, que ilustra a criação, execução e sincronização de duas \textit{threads}.
\begin{Verbatim}[numbers=left]
#include <stdio.h>
#include <pthread.h>
int x = 0, y = 0; // Variáveis compartilhadas
void funcao1(void *threadarg) {
    x = 1;
    ... // várias instruções
    if (y == 0) printf("1 ");
    pthread_exit(0);
}
void funcao2(void *threadarg) {
    y = 1;
    ... // várias instruções
    if (x == 0) printf("2 ");
    pthread_exit(0);
}
void main() {
    pthread_t t1, t2;
    pthread_create(&t1, NULL, (void *)funcao1, NULL);
    pthread_create(&t2, NULL, (void *)funcao2, NULL);
    pthread_join(t1, NULL);
    pthread_join(t2, NULL);
}
\end{Verbatim}
Ao final da execução da função \texttt{main}, será impresso
\begin{QuestionOptions}
    \item ambos os valores ``1'' e ``2''.
    \item o valor ``1'', necessariamente.
    \item o valor ``2'', necessariamente.
    \item o valor ``1'', ou o valor ``2'', mas nunca ambos.
    \item o valor ``1'', ou o valor ``2'', ou nenhum valor, mas nunca ambos.
\end{QuestionOptions}
\end{Question}

Esse código ilustra uma condição de corrida em memória compartilhada sem travas,
em que o resultado depende do entrelaçamento das instruções pelo escalonador.
Temos três casos sequenciais possíveis \parencite[seç.~2.3]{tanenbaum2016}:
\begin{itemize}
    \item A \texttt{funcao1} executa inteiramente primeiro.
    \texttt{x} recebe 1.
    Verifica \texttt{y}, que ainda é 0, e imprime ``1''.
    Em seguida, \texttt{funcao2} roda: \texttt{y} recebe 1.
    Verifica \texttt{x} (agora 1), então não imprime nada.
    \item A \texttt{funcao2} executa inteiramente primeiro.
    \texttt{y} recebe 1.
    Verifica \texttt{x} (ainda é 0), imprime ``2''.
    Em seguida, \texttt{funcao1} roda: \texttt{x} recebe 1.
    Verifica \texttt{y} (agora 1), e não imprime nada.
    \item A execução se entrelaça e ambas as threads chegam às suas primeiras linhas antes de continuar.
    A thread 1 executa \texttt{x = 1}.
    O SO tira a CPU da thread 1 e entrega à thread 2, que executa \texttt{y = 1}.
    A partir daí, qualquer ordem das avaliações dos comandos \texttt{if} resultará em falso,
    pois tanto \texttt{x} quanto \texttt{y} já são 1.
    O programa termina imprimindo nenhum valor.
\end{itemize}

O único cenário impossível nesse modelo é a impressão de ambos ``1'' e ``2'',
pois para que os dois \texttt{if} avaliem verdadeiro simultaneamente,
nenhuma das \textit{threads} poderia ter executado sua primeira instrução de atribuição ainda,
evidenciando que a resposta correta é a opção E.

\begin{Question}{2017-35}
Um programador inexperiente está desenvolvendo um sistema \textit{multithread} que possui duas estruturas
de dados diferentes, $E_1$ e $E_2$, as quais armazenam valores inteiros. O acesso concorrente a essas
estruturas é controlado por semáforos. Durante sua execução, o sistema dispara as \textit{threads} $T_1$ e $T_2$
simultaneamente. A tabela a seguir possibilita uma visão em linhas gerais dos algoritmos dessas \textit{threads}.
\begin{center}
\small
\begin{tabular}{|l|l|}\hline
\rowcolor{lightgray}\textbf{$T_1$} & \textbf{$T_2$} \\\hline
Aloca $E_1$ & Aloca $E_2$ \\\hline
Calcula a média $M_1$ dos valores de $E_1$ & Calcula a soma $S_1$ de todos os valores de $E_2$ \\\hline
Aloca $E_2$ & Aloca $E_1$ \\\hline
Calcula a média $M_2$ dos valores de $E_2$ & Calcula a soma $S_2$ de todos os valores de $E_1$ \\\hline
Calcula $M_3 = M_1 + M_2$ & Calcula $S_3 = \lvert S_1 - S_2 \rvert$ \\\hline
Soma $M_3$ em todos os valores de $E_2$ & Subtrai $S_2$ de todos os valores de $E_1$ \\\hline
Libera $E_1$ & Libera $E_2$ \\\hline
Libera $E_2$ & Libera $E_1$ \\\hline
\end{tabular}
\end{center}
Durante a execução do referido programa, é possível que
\begin{QuestionOptions}
    \item não ocorra \textit{deadlock},
    porque a sequência de alocação dos recursos impede naturalmente o problema.
    \item ocorra \textit{deadlock},
    que pode ser evitado se o programador tomar o cuidado de não executar cálculos entre um pedido de alocação e outro.
    \item ocorra \textit{deadlock},
    sendo a probabilidade dessa ocorrência tão baixa e sua consequência tão inócua que não haverá comprometimento do programa.
    \item não ocorra \textit{deadlock},
    desde que o programador use semáforos para controlar o acesso às estruturas de dados,
    o que é suficiente para evitar o problema.
    \item ocorra \textit{deadlock},
    que pode ser evitado se o programador tomar o cuidado de solicitar
    o acesso às estruturas de dados na mesma ordem em ambas as \textit{threads}.
\end{QuestionOptions}
\end{Question}

Analisando a sequência de alocação:
a \textit{thread} $T_1$ solicita o recurso $E_1$ e, em seguida, $E_2$;
já a \textit{thread} $T_2$ solicita $E_2$ e, posteriormente, $E_1$.
Se $T_1$ alocar $E_1$ no mesmo instante em que $T_2$ alocar $E_2$,
$T_1$ ficará bloqueada aguardando $E_2$, enquanto $T_2$ ficará bloqueada aguardando $E_1$,
caracterizando uma espera circular e resultando em um \textit{deadlock}.

Para prevenir esse problema, uma estratégia clássica de prevenção consiste em impor uma ordenação global na alocação de recursos.
Se ambas as \textit{threads} solicitarem o acesso às estruturas de dados na mesma ordem (por exemplo, alocando $E_1$ antes de $E_2$),
a condição de espera circular é eliminada \parencite[seç.~6.5]{tanenbaum2016}.
A opção correta é a E.

\begin{Question}{2021-17}
Durante parte do tempo, um processo está ocupado realizando computações internas e outras coisas que não levam a condições de corrida.
No entanto, às vezes, um processo tem de acessar uma memória compartilhada ou arquivos,
ou realizar outras tarefas críticas que podem levar a corridas.
Essa parte do programa onde a memória compartilhada é acessada é chamada de \emph{região crítica} ou \emph{seção crítica}.
Se conseguíssemos arranjar as coisas de maneira que dois processos jamais estivessem em suas regiões críticas ao mesmo tempo,
poderíamos evitar as corridas.
Embora essa exigência evite as condições de corrida,
ela não é suficiente para garantir que processos em paralelo cooperem de modo correto e eficiente usando dados compartilhados.
Precisamos que quatro condições se mantenham para chegar a uma boa solução.
\begin{enumerate}
    \item Dois processos jamais podem simultaneamente estar dentro de suas regiões críticas.
    \item Nenhuma suposição pode ser feita a respeito de velocidades ou de número de CPUs.
    \item Nenhum processo executando fora de sua região crítica pode bloquear qualquer processo.
    \item Nenhum processo deve ser obrigado a esperar eternamente para entrar em sua região crítica.
\end{enumerate}
\vspace{-0.5em}
\begin{flushright}
    \scriptsize
    TANENBAUM, A. S. \textbf{Sistemas Operacionais Modernos}.\\
    4. ed., São Paulo: Pearson, p. 83, 2016 (adaptado).
\end{flushright}
Considerando o texto e a figura [omitida] apresentados, avalie as asserções a seguir e a relação proposta entre elas.
\begin{enumerate}[label=\Roman*.]
    \item Em algumas situações, a exclusão mútua pode ser obtida por meio da desabilitação da interrupção controlada pelo Sistema Operacional,
    não sendo permitido que o seu controle seja feito pelo usuário.
\end{enumerate}
\begin{center}\textbf{PORQUE}\end{center}
\begin{enumerate}[label=\Roman*.,start=2]
    \item A desabilitação da interrupção é uma técnica que pode impedir que o processador que está executando um processo
    em sua região crítica seja interrompido para executar outro código,
    sendo mais eficiente em sistemas de multiprocessadores devido a quantidade de processos concorrentes.
\end{enumerate}
A respeito dessas asserções, assinale a opção correta.
\begin{QuestionOptions}
    \item As asserções I e II são proposições verdadeiras, e a II é uma justificativa correta da I.
    \item As asserções I e II são proposições verdadeiras, mas a II não é uma justificativa correta da I.
    \item A asserção I é uma proposição verdadeira, e a II é uma proposição falsa.
    \item A asserção I é uma proposição falsa, e a II é uma proposição verdadeira.
    \item As asserções I e II são proposições falsas.
\end{QuestionOptions}
\end{Question}

A asserção I é verdadeira.
O próprio sistema operacional pode desabilitar interrupções temporariamente ao atualizar estruturas internas do núcleo
para garantir exclusão mútua em um único processador;
tal poder nunca é concedido ao modo usuário, pois um processo malicioso poderia monopolizar a CPU indefinidamente.

A asserção II é falsa.
Desabilitar interrupções afeta apenas o núcleo/CPU que executou a instrução.
Em sistemas multiprocessados, outros núcleos continuam executando em paralelo e acessando a memória compartilhada normalmente,
tornando essa técnica ineficaz para garantir exclusão mútua entre múltiplos processadores \parencite[seç.~2.3.1--3]{tanenbaum2016}.
A opção correta é a C.

\printbibliography[heading=bibintoc]
\end{document}
